\chapter{Obiectivele cercetării}

Obiectivul principal al cercetării constă în evaularea tehnologiilor
moderne de comunicare din domeniul sistemelor distribuite, în scopul
implementării unui nou sistem de compilare distribuită pentru proiecte C/C++, care
să ofere performanțe competitive în ceea ce privește eficientizarea
procesului de compilare, comparativ cu alternativele deja existente.
Sistemul dezvoltat propune o abordare eficientă, scalabilă și flexibilă
din punct de vedere arhitectural, de a crește viteza procesului de 
compilare, fără a introduce incompatibilități cu uneltele deja folosite
în dezvoltare proiectelor.

Compilarea proiectelor dezvoltate folosind C/C++ este compusă din etape
multiple de transformare a codului sursă: preprocesarea, compilarea,
generarea de cod mașină pentru platforma vizată și link-editarea. Unele 
operațiuni aferente procesului de compilare pot fi executate independent,
în paralel, însă gestiunea dependințelor dintre fișierele componente
ale proiectului, artefactele generate și setările de configurare ale compilatorului 
introduce un grad însemnat de complexitate. Un obiectiv ce trebuie 
îndeplinit de către orice sistem distribuit de acest tip îl reprezintă
identificarea corectă a sarcinilor care pot fi preluate de către sistemele
de calcul auxiliare, spre a fi paralelizate, astfel încât să fie menținută
corectitudinea și capacitatea de reproducere facilă a procesului de compilare.

Lucrarea investighează modul în care sarcinile de compilare pot să
fie distribuite într-o manieră eficientă către mai multe noduri de calcul din
rețea, pentru a se minimiza suprautilizarea resurselor de calcul, precum și a
canalelor de comunicare. De asemenea, sistemul trebuie să păstreze în evidență
dsiponibilitatea sistemelor ce contribuie la procesul de build și să își mențină
o stare de funcționare consistentă în cazul apariției erorilor în timpul 
transmisiei de date în rețea.

Astfel, metodologia de întocmire a lucrării de cercetare propune realizarea
unei serii de mai multe activități, pentru a explora modul de funcționare și
utilitatea conceptului de compilare distribuită pentru proiectede dezvoltate
folosind limbajele de programare \textbf{C} și \textbf{C++}:

\newpage

\begin{itemize}
    \item Studiul bibliografic - în vederea dezvoltării unui sistem de compilare
    distribuită, este necesară o analiză aprofundată a soluțiilor deja existente
    pentru realizarea acestui tip de sarcină. Studiul bibliografic oferă o 
    imagine de ansamblu a principiilor ce stau la baza compilării
    distribuite, a implicațiilor practice ale aplicării unei astfel de abordări,
    precum și a cazurilor de utilizare cele mai frecvente pentru astfel de sisteme
    \item Proiectarea arhitecturii sistemului - aceasta reprezintă o etapă premergătoare
    implementării concrete și cuprinde descrierea comportamentului pentru modulele
    care intră în componența sistemului, cu modalitățile de funcționare și interacțiune
    aferente. De asemenea, se studiază posibilitatea aplicării unor concepte moderne
    pentru dezvoltarea sistemelor distribuite și se definește structura arhitecturii
    hardware și software a ansamblului
    \item Implementarea sistemului - pe baza arhitecturii definite în etapa anterioară,
    se implementeză soluții practice pentru îndeplinirea operațiunilor asociate funcționării
    sistemului de compilare distribuită: prelucrarea datelor de intrare, balansarea sarcinii
    de compilare între nodurile participante, stocarea artefactelor rezultate, gestiunea
    comunicării dintre noduri pe baza unui protocol predefinit
    \item Crearea unui mediu de test - validarea funcționării corecte a sistemului
    dezvoltate poate avea loc doar în contextul utilizării unei configurații
    corespunzătoare, care să fie cât mai apropiată de condițiile de operare reale, în
    ceea ce privește dispozitivele hardware implicate, atât din punct de vedere al puterii
    de calcul, dar și al resurselor de rețea
    \item Analiza rezultatelor obținute - pe baza metricilor ce caracterizează 
    performanțele și comportamentul sistemului se pot identifica principalele avantaje și 
    dezavantaje ale abordării alese, raportate la celelalte alternative propuse în 
    literatura de specialitate, dar și scenariile optime de utilizare pentru distribuirea 
    procesului de compilare folosind mai multe sisteme de calcul 
\end{itemize}

Pentru a restrânge aria de cercetare asupra conceptului de compilare distribuită
pentru \textbf{C/C++}, nu se vor studia amănunțit toate aspectele adiacente funcționării
unui astfel de sistem, precum testarea mai multor configurații de rețea, asigurarea 
funcționării compatibilității aplicației cu platform de rulare multiple sau optimizarea
procesului de deployment.

Se urmărește dezvoltarea unui sistem distribuit care să fie compatibil cu sistemele de operare
bazate pe \textbf{Linux} și uneltele software uzuale în dezvoltarea aplicațiilor C/C++, dintre
care amintim: compilatoarele GCC și Clang, respectiv sistemele de build CMake, Ninja și GNU Make.
Astfel, ne limităm doar la asigurarea interacțiunii facile a sistemului cu  cele mai cunoscute
medii de dezvoltare moderne.

De asemenea, funcționarea sistemul distribuit presupune apartenența sistemelor de calcul
participante la o rețea de calculatoare locală LAN. Scenariul funcționării într-o rețea
de noduri distribuite asupra unei arii geografice însemnate implică necesitatea implementării
unor măsuri de securitate suplimentare pentru păstrarea integrității și confidențialității
datelor transferate, iar această activitate depășește sfera de cercetare a lucrării.

Totodată, o altă precondiție care este impusă pentru funcționarea o constituie uniformitatea
mediului de rulare pentru nodurile participante, având la bază presupunerea că sistemele de calcul
implicate sunt dotate cu aceeași versiune de sistem de operare, sisteme de build și compilatoare, 
iar arhitecturile hardware ale acestora sunt identice. Nu se studiază metode de partiționare 
ale colecției de noduri de calcul pe baza compatibilității dintre platformele de rulare ale acestora.

După implementarea sistemului de compilare distribuită, un alt obiectiv important al lucrării îl constituie
evaluarea procesului de utilizare a sistemului, având în vedere analiza următoarelor aspecte cheie:

\begin{itemize}
    \item Durata de build - se măsoară timpul de compilare pentru mai multe proiecte open-source arhicunoscute,
    atât folosind un unic nod de compilare, pentru a se determina măsura în care utilizarea sistemului
    introduce overhead, raportat la durata de compilare locală pe nodul respectiv, precum și folosind
    multiple noduri care contribuie la procesul de compilare, pentru a măsura, în mod evident, îmbunătățirile
    aduse la eficientizarea procesului de build
    \item Comportamenul determinist - în mod ideal, comportamnetul procesului de build, folosind sistemul
    dezvoltat, trebuie să fie același la fiecare rulare, iar erorile sau warning-urile generate de către
    compilator trebuie să fie cât mai apropiate de cele rezultate în urma procesului de compilare locală
    \item Randamentul paralelizării - în funcție de natura proiectului care este supus procesului de compilare,
    sistemul de compilare distribuită aduce grade diferite de îmbunătățire a timpilor de build. În scenariul ideal,
    structura proiectului permite o micșorare a timpului de compilare direct proporțională cu numărul de sarcini de
    compilare paralelă oferite de totalitatea nodurilor participante, însă în realitate, unele proiecte nu se bucură
    de îmbunătățiri semnificative folosind o astfel de abordare
    \item Volumul de date transferat - distribuirea procesului de compilare către mai multe sisteme de calcul poate
    să implice transferul unor volume de date semnificative, în rețea, ceea ce introduce necesitatea accesului la 
    resurse de rețea adecvate, pentru a nu deteriora eficiența procesului de build din cauza carențelor date de 
    capacitatea de transfer dintr-o perioadă de timp sau latența de inițiere a transferului
\end{itemize}